% Options for packages loaded elsewhere
\PassOptionsToPackage{unicode}{hyperref}
\PassOptionsToPackage{hyphens}{url}
\PassOptionsToPackage{dvipsnames,svgnames,x11names}{xcolor}
%
\documentclass[
  10pt,
]{article}
\usepackage{amsmath,amssymb}
\usepackage{iftex}
\ifPDFTeX
  \usepackage[T1]{fontenc}
  \usepackage[utf8]{inputenc}
  \usepackage{textcomp} % provide euro and other symbols
\else % if luatex or xetex
  \usepackage{unicode-math} % this also loads fontspec
  \defaultfontfeatures{Scale=MatchLowercase}
  \defaultfontfeatures[\rmfamily]{Ligatures=TeX,Scale=1}
\fi
\usepackage{lmodern}
\ifPDFTeX\else
  % xetex/luatex font selection
  \setmainfont[]{DejaVu Serif}
  \setsansfont[]{DejaVu Sans}
  \setmonofont[]{DejaVu Sans Mono}
  \setmathfont[]{Latin Modern Math}
\fi
% Use upquote if available, for straight quotes in verbatim environments
\IfFileExists{upquote.sty}{\usepackage{upquote}}{}
\IfFileExists{microtype.sty}{% use microtype if available
  \usepackage[]{microtype}
  \UseMicrotypeSet[protrusion]{basicmath} % disable protrusion for tt fonts
}{}
\makeatletter
\@ifundefined{KOMAClassName}{% if non-KOMA class
  \IfFileExists{parskip.sty}{%
    \usepackage{parskip}
  }{% else
    \setlength{\parindent}{0pt}
    \setlength{\parskip}{6pt plus 2pt minus 1pt}}
}{% if KOMA class
  \KOMAoptions{parskip=half}}
\makeatother
\usepackage{xcolor}
\usepackage[margin=2.0cm]{geometry}
\usepackage{longtable,booktabs,array}
\usepackage{calc} % for calculating minipage widths
% Correct order of tables after \paragraph or \subparagraph
\usepackage{etoolbox}
\makeatletter
\patchcmd\longtable{\par}{\if@noskipsec\mbox{}\fi\par}{}{}
\makeatother
% Allow footnotes in longtable head/foot
\IfFileExists{footnotehyper.sty}{\usepackage{footnotehyper}}{\usepackage{footnote}}
\makesavenoteenv{longtable}
\setlength{\emergencystretch}{3em} % prevent overfull lines
\providecommand{\tightlist}{%
  \setlength{\itemsep}{0pt}\setlength{\parskip}{0pt}}
\setcounter{secnumdepth}{-\maxdimen} % remove section numbering
% doc_header.tex -- LaTeX preamble additions used by build_docs.sh for every
% Markdown -> LaTeX -> PDF conversion (supplements, README, submission files).
\usepackage{etoolbox}
% longtable may not be loaded by pandoc when a document has no table
\usepackage{longtable}
% bold vectors: \bm -> \symbf under unicode-math (XeLaTeX), bm package otherwise
\makeatletter
\@ifpackageloaded{unicode-math}{\providecommand{\bm}[1]{\symbf{#1}}}{\usepackage{bm}}
\makeatother
% smaller type in long tables so that wide audit tables fit the text width;
% inside a table a line break is allowed after an underscore of plain text
% (pandoc writes it as \_), so that long identifiers wrap inside their column
% rows are set a little apart, so that symbols with indices do not touch
\AtBeginEnvironment{longtable}{\footnotesize\renewcommand{\arraystretch}{1.45}%
  \let\sptunderscore\_\def\_{\sptunderscore\allowbreak}}
\setlength{\LTcapwidth}{\textwidth}
% allow line breaks after slashes and hyphens inside table cells
\usepackage{xurl}
% code blocks: break over-long lines instead of running into the margin.
% Plain ``` blocks are typeset by pandoc as {verbatim}, highlighted ones
% (```bash) as {Highlighting}; both are re-declared as breakable fancyvrb
% environments.  The continuation line is marked with a small hooked arrow.
\usepackage{fvextra}
\RecustomVerbatimEnvironment{verbatim}{Verbatim}{breaklines,breakanywhere,fontsize=\small}
\makeatletter
\@ifundefined{Highlighting}{}{%
  \RecustomVerbatimEnvironment{Highlighting}{Verbatim}{breaklines,breakanywhere,fontsize=\small,commandchars=\\\{\}}}
\makeatother
% inline code (pandoc writes it as \texttt{...}): never hyphenate -- an inserted
% hyphen could be read as part of a file name -- but allow a line break after
% an underscore, slash, comma, colon or equals sign, so that long paths and
% column lists wrap inside the text width and inside table columns.
% Inside a formula (e.g. the name of a notebook symbol in an equation) the same
% command must first leave math mode, which \text of amsmath does.
\ExplSyntaxOn
\tl_new:N \l__spt_code_tl
\cs_new_protected:Npn \__spt_texttt:n #1
  {
    \tl_set:Nn \l__spt_code_tl {#1}
    \regex_replace_all:nnN { ([\/\,\:\=]) } { \1 \c{allowbreak} } \l__spt_code_tl
    \regex_replace_all:nnN { \c{_} } { \c{_} \c{allowbreak} } \l__spt_code_tl
    \group_begin: \ttfamily \hyphenchar \font = -1 \scan_stop: \l__spt_code_tl \group_end:
  }
\RenewDocumentCommand \texttt { m }
  {
    \mode_if_math:TF { \text { \__spt_texttt:n {#1} } } { \__spt_texttt:n {#1} }
  }
\ExplSyntaxOff
% A figure is written in the Markdown source as an image paragraph followed by
% a caption paragraph ("**Figure S1.** ..."), so that the caption is also
% visible on GitHub.  An image that starts a paragraph is centered.
\usepackage{graphicx}% pandoc loads it only for documents that contain an image
\NewCommandCopy\sptincludegraphics\includegraphics
\RenewDocumentCommand\includegraphics{O{} m}{%
  \ifvmode\begin{center}\sptincludegraphics[#1]{#2}\end{center}%
  \else\sptincludegraphics[#1]{#2}\fi}
% keep figures where they are referenced
\usepackage{float}
\floatplacement{figure}{H}
% running header with the document title
\usepackage{fancyhdr}
\pagestyle{fancy}
\fancyhf{}
\fancyhead[L]{\small PROJECT spacetimes --- documentation and provenance}
\fancyhead[R]{\small\thepage}
\renewcommand{\headrulewidth}{0.4pt}
% 16x16 matrices: tighter columns so that a full matrix fits the text width
\setlength{\arraycolsep}{3.2pt}
% The documents of this project contain long tagged symbols, file paths and
% identifiers that cannot be hyphenated.  A paragraph is therefore allowed to
% have looser lines before TeX gives up and lets a line stick out of the margin.
\AtBeginDocument{\tolerance=2000 \emergencystretch=5em \hbadness=2000}
\ifLuaTeX
  \usepackage{selnolig}  % disable illegal ligatures
\fi
\IfFileExists{bookmark.sty}{\usepackage{bookmark}}{\usepackage{hyperref}}
\IfFileExists{xurl.sty}{\usepackage{xurl}}{} % add URL line breaks if available
\urlstyle{same}
\hypersetup{
  pdftitle={P12 --- The sixteen-component spinor representation of Spin(4,4) and of Spin(7,1), and its two chiral parts},
  pdfauthor={Prepared for Patrick Lee Nash, Ph.D.},
  colorlinks=true,
  linkcolor={blue},
  filecolor={Maroon},
  citecolor={blue},
  urlcolor={blue},
  pdfcreator={LaTeX via pandoc}}

\title{P12 --- The sixteen-component spinor representation of Spin(4,4)
and of Spin(7,1), and its two chiral parts}
\usepackage{etoolbox}
\makeatletter
\providecommand{\subtitle}[1]{% add subtitle to \maketitle
  \apptocmd{\@title}{\par {\large #1 \par}}{}{}
}
\makeatother
\subtitle{Provenance of PROJECT spacetimes: check and correction of two
claims of the author}
\author{Prepared for Patrick Lee Nash, Ph.D.}
\date{November 15, 2026}

\begin{document}
\maketitle

\hypertarget{the-two-claims-and-the-verdict-in-brief}{%
\section{1. The two claims, and the verdict in
brief}\label{the-two-claims-and-the-verdict-in-brief}}

On November 15, 2026, the author asked for the following two claims to
be checked, verified and corrected in a provenance:

\begin{quote}
Understand and verify that the 16 by 16 dimensional representation of
Spin(4,4) is reducible to direct sum of two 8 by 8 dimensional
irreducible representations of Spin(4,4)\_left and Spin(4,4)\_right.

Understand and verify that the 16 by 16 dimensional representation of
Spin(7,1) is reducible to direct sum of 8 by 8 dimensional irreducible
representations of Spin(7,1)\_left and Spin(7,1)\_right.
\end{quote}

\textbf{Claim 1 (4+4) is correct in substance and needs two corrections
of wording.} The real sixteen-dimensional spinor representation of
\(\mathrm{Spin}_0(4,4)\) is the direct sum of two \emph{real}
eight-dimensional invariant subspaces, the left-handed and the
right-handed spinors. Each is irreducible, over the real and over the
complex numbers, and the two are not equivalent. Corrections: (a)
``Spin(4,4)\_left'' and ``Spin(4,4)\_right'' are not two groups; they
are two representations of the \emph{one} group
\(\mathrm{Spin}_0(4,4)\), the left and right half-spin representations.
(b) Neither half-spin representation is faithful: each sends two
elements of the group to the identity, so that the group of \(8\times8\)
matrices that each produces is not a copy of \(\mathrm{Spin}_0(4,4)\)
but a quotient of it by a group of order two, which preserves a
symmetric form of signature (4,4).

\textbf{Claim 2 (7+1) is correct only over the complex numbers.} The
sixteen-dimensional representation of \(\mathrm{Spin}_0(7,1)\) that the
project uses is real, and as a real representation it is
\emph{irreducible}: no real change of basis splits it. Over the complex
numbers it is the direct sum of two eight-dimensional irreducible
representations, the left-handed and the right-handed spinors, which are
not equivalent; the right one is the complex conjugate of the left one,
and also its dual. Each is faithful. The same correction (a) applies:
the two are representations of one group. The \(8\times8\) matrices of
the two parts are necessarily complex.

The rest of this document defines every term, gives the method, the
results with the check that verifies each, the comparison with the
four-dimensional case of E. A. Lord, and the complete output of the
program.

\hypertarget{definitions}{%
\section{2. Definitions}\label{definitions}}

\emph{The Dirac matrices.} \(\gamma^A\), \(A=0,\dots,7\), are the real
\(16\times16\) matrices of the project, built by the program from the
author's notebooks, with
\[\gamma^A\gamma^B+\gamma^B\gamma^A=2\eta^{AB},\]
\[\eta^{(4+4)}=\mathrm{diag}(+1,+1,+1,+1,-1,-1,-1,-1),\]
\[\eta^{(7+1)}=\mathrm{diag}(-1,-1,-1,-1,+1,-1,-1,-1).\] The coordinates
are \(x0,\dots,x7\): \(x1,x2,x3\) ordinary space, \(x4\) ordinary time,
\(x0\) the hidden coordinate, \(x5,x6,x7\) the extra coordinates.

\emph{The spin group.} The 28 matrices
\[\Sigma^{AB}=\tfrac14\,[\gamma^A,\gamma^B],\qquad A<B,\] span a Lie
algebra isomorphic to \(\mathfrak{so}(4,4)\) or \(\mathfrak{so}(7,1)\),
the Lie algebra of the rotations of the spacetime:
\([\Sigma^{AB},\gamma^C]=\eta^{BC}\gamma^A-\eta^{AC}\gamma^B\). The
group generated by the exponentials
\(\exp(\sum_{A<B}w_{AB}\Sigma^{AB})\) with real numbers \(w_{AB}\) is
\(\mathrm{Spin}_0(4,4)\) or \(\mathrm{Spin}_0(7,1)\), the identity
component of the spin group; the subscript 0 says that only elements
connected to the identity are meant. Acting on the sixteen components of
a spinor, these matrices form the \textbf{spinor representation}, of
dimension sixteen.

\emph{Chirality.} \(\omega=\gamma^{0}\gamma^{1}\cdots\gamma^{7}\); the
chirality matrix is \(\chi=\omega\) on 4+4 and \(\chi=i\omega\) on 7+1;
the chiral projectors are \(P^L=\tfrac12(1-\chi)\) and
\(P^R=\tfrac12(1+\chi)\). A spinor with \(P^L\psi=\psi\) is
\textbf{left-handed}, one with \(P^R\psi=\psi\) \textbf{right-handed}.

\emph{The charge-conjugation matrix.}
\(C^-=\gamma^{0}\gamma^{1}\gamma^{2}\gamma^{3}\) on 4+4 and
\(C^-=\gamma^{4}\omega\) on 7+1; it defines the bilinear form
\(\psi^TC^-\psi'\) of two spinors.

\emph{Terms of representation theory.} A subspace is \textbf{invariant}
if every matrix of the representation maps it into itself. A
representation is \textbf{irreducible} if it has no invariant subspace
other than zero and the whole space, and \textbf{reducible} otherwise;
whether it is irreducible can depend on whether real or complex numbers
are allowed as coefficients. Two representations are \textbf{equivalent}
if an invertible matrix \(T\) turns the matrices of one into those of
the other, \(T\,\rho_1(g)=\rho_2(g)\,T\) for every group element \(g\).
A representation is \textbf{faithful} if only the identity element is
represented by the unit matrix. The \textbf{commutant} of a
representation is the set of matrices that commute with all its
matrices; for the representation of a connected group it is the set of
matrices that commute with all the generators \(\Sigma^{AB}\).

\hypertarget{the-method}{%
\section{3. The method}\label{the-method}}

Two theorems reduce the questions to the commutant.

\begin{itemize}
\tightlist
\item
  \emph{Complete reducibility} (Weyl). The Lie algebras
  \(\mathfrak{so}(4,4)\) and \(\mathfrak{so}(7,1)\) are semisimple.
  Every finite-dimensional representation of a semisimple Lie algebra,
  and of its connected group, is completely reducible: every invariant
  subspace has an invariant complement. An invariant subspace with an
  invariant complement is the image of a projector (a matrix \(E\) with
  \(E^2=E\)) in the commutant.
\item
  \emph{Schur's lemma.} Over the complex numbers, a representation that
  is a direct sum of irreducible parts, the \(k\)-th of which occurs
  \(n_k\) times, has a commutant of dimension \(\sum_kn_k^2\).
\end{itemize}

Hence: over the real numbers the representation is irreducible exactly
when its real commutant contains no projector other than 0 and 1. Over
the complex numbers, a commutant of dimension 2, together with two
non-zero invariant subspaces, means exactly two irreducible parts, each
occurring once, and therefore not equivalent.

The commutant is the solution space of the homogeneous linear system
\(XM^{AB}=M^{AB}X\) for the 256 entries of a \(16\times16\) matrix
\(X\), with \(M^{AB}=4\Sigma^{AB}=[\gamma^A,\gamma^B]\), whose entries
are integers. The program computes the rank of this system modulo a
large prime. That number is a lower bound for the rank over the rational
numbers, so that \(256-\text{rank}\) is an upper bound for the dimension
of the commutant; the program then exhibits that many independent
solutions, which proves the dimension. Because the coefficients are
real, the complex commutant has the same dimension as the real one.
Every other check is an identity between integer (or, on 7+1,
Gaussian-integer) matrices, and is exact.

\hypertarget{results}{%
\section{4. Results}\label{results}}

\hypertarget{the-44-spacetime}{%
\subsection{4.1 The 4+4 spacetime}\label{the-44-spacetime}}

\begin{enumerate}
\def\labelenumi{\arabic{enumi}.}
\tightlist
\item
  The matrices satisfy the Clifford relations with \(\eta^{(4+4)}\)
  (check \texttt{44-R0}).
\item
  \(\omega\) is real, \(\omega^2=+1\), \(\mathrm{Tr}\,\omega=0\); it
  anticommutes with every \(\gamma^A\) and therefore commutes with every
  \(\Sigma^{AB}\) (\texttt{44-R1}).
\item
  The commutant of the 28 generators is spanned by \(1\) and \(\omega\)
  and has dimension 2 (\texttt{44-R2}; rank 254 of the linear system).
\item
  \(P^L=\tfrac12(1-\omega)\) and \(P^R=\tfrac12(1+\omega)\) are
  \textbf{real} projectors of rank 8 in the commutant: the real
  representation splits into two real eight-dimensional invariant
  subspaces. Since the commutant is spanned by \(P^L\) and \(P^R\), each
  part is irreducible over the complex numbers (hence also over the real
  numbers), and the two parts are not equivalent (\texttt{44-R3}).
\item
  In the basis of the matrix files the first eight components are
  left-handed and the last eight right-handed:
  \(\omega=\mathrm{diag}(-1_8,+1_8)\), and every \(\Sigma^{AB}\) is
  block-diagonal (\texttt{44-R4}). With the blocks
  \(\gamma^A=\begin{pmatrix}0&\bar\tau^A\\ \tau^A&0\end{pmatrix}\) of
  the files, multiplying the block matrices shows that the upper-left
  (left-handed) block of \(\Sigma^{AB}\) is
  \(\tfrac14(\bar\tau^A\tau^B-\bar\tau^B\tau^A)\) and the lower-right
  (right-handed) block is
  \(\tfrac14(\tau^A\bar\tau^B-\tau^B\bar\tau^A)\).
\item
  The 28 generators remain linearly independent on each part: the whole
  Lie algebra acts faithfully on each part (\texttt{44-R6}).
\item
  The form \(\psi^TC^-\psi'\) is invariant, \(C^-\) is symmetric with
  \((C^-)^2=1\), and it commutes with \(\chi\): it pairs each part with
  itself, and on each part it has the eigenvalues \(+1\) and \(-1\) four
  times each. Each part therefore carries an invariant symmetric form of
  signature (4,4), and the matrices of each part lie in the group
  \(SO(4,4)\) of that form (\texttt{44-R7}).
\item
  \(\omega=(\gamma^{0}\gamma^{1})(\gamma^{2}\gamma^{3})(\gamma^{4}\gamma^{5})(\gamma^{6}\gamma^{7})\),
  and each factor squares to \(-1\), so that each factor equals
  \(\exp(\tfrac\pi2\gamma^a\gamma^b)=\exp(\pi\Sigma^{ab})\): \(\omega\)
  and \(-\omega\) belong to \(\mathrm{Spin}_0(4,4)\). On the left part
  \(\omega=-1\); hence the element \(-\omega\) is represented on the
  left part by the unit matrix, and the element \(\omega\) on the right
  part (\texttt{44-R8}). \textbf{Neither part is a faithful
  representation of the group}; the sixteen-dimensional sum is. The
  group of \(8\times8\) matrices produced by the left part is
  \(\mathrm{Spin}_0(4,4)/\{1,-\omega\}\), that of the right part
  \(\mathrm{Spin}_0(4,4)/\{1,\omega\}\). Each has the dimension 28 of
  \(SO(4,4)\) and lies in it (item 7). (That each is all of the identity
  component of \(SO(4,4)\), and that \(\pm1,\pm\omega\) are all the
  central elements, are known facts, the ``triality'' of
  eight-dimensional rotations; they were not computed.)
\end{enumerate}

\hypertarget{the-71-spacetime}{%
\subsection{4.2 The 7+1 spacetime}\label{the-71-spacetime}}

\begin{enumerate}
\def\labelenumi{\arabic{enumi}.}
\tightlist
\item
  The Clifford relations hold with \(\eta^{(7+1)}\) (\texttt{71-R0}).
\item
  \(\omega\) is real, \(\omega^2=-1\), \(\mathrm{Tr}\,\omega=0\); it
  commutes with every \(\Sigma^{AB}\) (\texttt{71-R1}).
\item
  The commutant is spanned by \(1\) and \(\omega\), of dimension 2
  (\texttt{71-R2}).
\item
  Because \(\omega^2=-1\), the real commutant \(\{a+b\,\omega\}\) is a
  copy of the complex numbers: \((a+b\omega)^2=a+b\omega\) has only the
  solutions 0 and 1. It contains no projector other than 0 and 1, and
  \textbf{the real sixteen-dimensional representation is irreducible
  over the real numbers} (\texttt{71-R3}).
\item
  Over the complex numbers, \(P^L=\tfrac12(1-i\omega)\) and
  \(P^R=\tfrac12(1+i\omega)\) split it into two eight-dimensional
  invariant subspaces, which by item 3 are irreducible and not
  equivalent (\texttt{71-R3}). In the basis of the files the generators
  are not block-diagonal (\texttt{71-R4}).
\item
  Because \(\omega\) and the generators are real, complex conjugation
  maps the left part onto the right part and commutes with the group:
  \textbf{the right part is the complex conjugate of the left part}
  (\texttt{71-R5}).
\item
  The generators remain independent on each part (\texttt{71-R6}).
\item
  The invariant form \(\psi^TC^-\psi'\) pairs the left part with the
  right part, because \(C^-\) anticommutes with \(\chi\): the right part
  is also the \textbf{dual} of the left part (\texttt{71-R7}).
\item
  \(\omega\) anticommutes with every \(\gamma^A\), so that it represents
  the inversion of all eight axes; that inversion reverses the direction
  of time, and \(\omega\) is therefore not in \(\mathrm{Spin}_0(7,1)\).
  The only central elements of \(\mathrm{Spin}_0(7,1)\) are \(+1\) and
  \(-1\), and \(-1\) is represented by \(-1\) on each part: \textbf{each
  part is faithful} (\texttt{71-R8}; the statement about the central
  elements is a known fact, not computed).
\end{enumerate}

\hypertarget{comparison-with-the-four-dimensional-case-lord-1971}{%
\subsection{4.3 Comparison with the four-dimensional case (Lord
1971)}\label{comparison-with-the-four-dimensional-case-lord-1971}}

E. A. Lord, \emph{General relativity from gauge invariance}, Proc. Camb.
Phil. Soc. 69, 423 (1971), works on the four-dimensional spacetime of
signature (1,3). There the Dirac spinor belongs to the direct sum of two
two-dimensional representations of \(SL(2,C)\), the double cover of the
Lorentz group: one in which a spinor transforms with a unimodular
\(2\times2\) matrix, and one in which it transforms with the inverse of
the Hermitian conjugate of that matrix. In the Weyl representation the
\(4\times4\) transformation is block-diagonal with these two blocks; the
two representations are not equivalent, and each is faithful. The 7+1
case is the analogue: two complex parts, conjugate to each other, both
faithful. The 4+4 case differs: the two parts are real and dual to
themselves, and neither is faithful.

\hypertarget{the-corrected-statements}{%
\section{5. The corrected statements}\label{the-corrected-statements}}

\begin{quote}
\textbf{Claim 1, corrected.} The sixteen-dimensional real spinor
representation of \(\mathrm{Spin}_0(4,4)\) is the direct sum of two real
eight-dimensional representations of the same group, the left-handed
(\(\chi=-1\)) and the right-handed (\(\chi=+1\)) half-spin
representations. Each is irreducible over the real and over the complex
numbers, the two are not equivalent, and each preserves a symmetric
bilinear form of signature (4,4). Neither is faithful: the left one
represents \(-\omega\), the right one \(\omega\), by the unit matrix, so
that the matrix groups they produce are
\(\mathrm{Spin}_0(4,4)/\{1,-\omega\}\) and
\(\mathrm{Spin}_0(4,4)/\{1,\omega\}\). In the basis of the project's
matrix files the left part is the first eight components and the right
part the last eight.

\textbf{Claim 2, corrected.} The sixteen-dimensional real spinor
representation of \(\mathrm{Spin}_0(7,1)\) is irreducible over the real
numbers. Over the complex numbers it is the direct sum of two complex
eight-dimensional representations of the same group, the left-handed and
the right-handed ones, \(P^{L,R}=\tfrac12(1\mp i\omega)\). Each is
irreducible and faithful, they are not equivalent, and each is the
complex conjugate and the dual of the other.

\textbf{For both.} ``Spin\_left'' and ``Spin\_right'' name two
representations of one group, not two groups. A theory in which the
left-handed and the right-handed parts transform with \emph{independent}
group elements, for example a gauge theory with two independent gauge
fields, uses a larger group, the product of the two matrix groups, of
which the spin group is the diagonal; that is a choice of the model, not
a consequence of the decomposition.
\end{quote}

\hypertarget{what-was-computed-how-to-reproduce-it-and-what-was-not-computed}{%
\section{6. What was computed, how to reproduce it, and what was not
computed}\label{what-was-computed-how-to-reproduce-it-and-what-was-not-computed}}

The command \texttt{spinreps} of the program of the main project (file
\texttt{code/spacetimes-rs/src/spinrep\_cmd.rs}) makes the checks
printed in Sec. 7, all in exact integer arithmetic, and writes
\texttt{data/spinrep\_checks.txt} and
\texttt{data/spinrep\_summary.csv}. From the root of the repository:

\begin{verbatim}
cargo run --release --manifest-path code/spacetimes-rs/Cargo.toml -- spinreps
\end{verbatim}

The command ends with ``spinreps: N of N checks passed''. This document
is built with

\begin{verbatim}
python3 code/scripts/assemble_docs.py docs_src/P12_spin_representations.src.md
bash code/scripts/build_docs.sh provenance/P12_spin_representations.md
\end{verbatim}

\emph{Not computed:} that \(\pm1,\pm\omega\) are all the central
elements of \(\mathrm{Spin}_0(4,4)\) and \(\pm1\) all those of
\(\mathrm{Spin}_0(7,1)\); that the matrix group of each 4+4 part is all
of the identity component of \(SO(4,4)\); the components of the spin
groups other than the identity component. The documentation of the
program (\texttt{provenance/P8\_code}) does not yet list the command
\texttt{spinreps}.

\hypertarget{the-complete-output-of-the-checks}{%
\section{7. The complete output of the
checks}\label{the-complete-output-of-the-checks}}

\begin{verbatim}
# PROJECT spacetimes: the spinor representation of Spin(4,4) and Spin(7,1) on sixteen components and its two eight-dimensional parts
# generated by: spacetimes spinreps (code/spacetimes-rs/src/spinrep_cmd.rs); exact integer arithmetic
# Sigma^{AB} = (1/4)[gamma^A, gamma^B]; omega = gamma^0...gamma^7; chi = omega (4+4), i omega (7+1); P^L = (1 - chi)/2, P^R = (1 + chi)/2
[PASS] 44-R0  4+4: the sixteen real matrices gamma^A satisfy gamma^A gamma^B + gamma^B gamma^A = 2 eta^{AB} with eta = diag(+1, +1, +1, +1, -1, -1, -1, -1)
[PASS] 44-R1  4+4: omega = gamma^0 ... gamma^7 is real, omega^2 = +1, Tr omega = 0; omega anticommutes with every gamma^A and commutes with all 28 generators Sigma^{AB}, so that the eigenspaces of chi are invariant under Spin_0
[PASS] 44-R2  4+4: the linear system X M^{AB} = M^{AB} X (all 28 generators; 6656 integer equations for the 256 entries of X) has rank 254 modulo a prime, so its solutions form a space of dimension at most 2; 1 and omega are two independent solutions; hence the commutant is exactly the span of 1 and omega, of dimension 2, over the real numbers and (the equations being real) also over the complex numbers
[PASS] 44-R3  4+4: P^L = (1 - omega)/2 and P^R = (1 + omega)/2 are REAL projectors of rank 8 that commute with every generator: the real 16-dimensional representation is the direct sum of two real 8-dimensional invariant subspaces, left (chi = -1) and right (chi = +1). The commutant is spanned by P^L and P^R (check R2): by Schur's lemma each part is irreducible over the complex numbers, and hence also over the real numbers, and the two parts are not equivalent (no non-zero invariant map from one to the other)
[PASS] 44-R4  4+4: in the basis of the matrix files the generators are block-diagonal with two 8x8 blocks, and omega = diag(-1_8, +1_8): the first eight components are the left-handed part and the last eight the right-handed part
[PASS] 44-R6  4+4: the 28 generators remain linearly independent on each part (rank 28 on the left part, 28 on the right part): the whole Lie algebra so(4,4) acts faithfully on each 8-dimensional part; each part is a representation of the one group Spin_0(4,4), not of a subgroup
[PASS] 44-R7  4+4: C^- is real and symmetric, (C^-)^2 = 1, C^- M^{AB} + (M^{AB})^T C^- = 0 for all generators, and C^- commutes with omega: the form psi^T C^- psi' is invariant and pairs each part with itself; on each part it has the eigenvalues +1 and -1 four times each (2 Tr(P^L C^-) = 0, 2 Tr(P^R C^-) = 0), so that each part carries an invariant symmetric form of signature (4,4), and the image of Spin_0(4,4) in each part lies in SO(4,4)
[PASS] 44-R8  4+4: omega = (gamma^0 gamma^1)(gamma^2 gamma^3)(gamma^4 gamma^5)(gamma^6 gamma^7), and each factor squares to -1, so that each factor equals exp((pi/2) gamma^a gamma^b) = exp(pi Sigma^{ab}): omega and -omega lie in Spin_0(4,4). omega acts as -1 on the left part and as +1 on the right part. Hence the element -omega acts as the identity on the left part and the element omega as the identity on the right part: NEITHER 8-dimensional part is a faithful representation of Spin_0(4,4); the 16-dimensional sum is
[PASS] 71-R0  7+1: the sixteen real matrices gamma^A satisfy gamma^A gamma^B + gamma^B gamma^A = 2 eta^{AB} with eta = diag(-1, -1, -1, -1, +1, -1, -1, -1)
[PASS] 71-R1  7+1: omega = gamma^0 ... gamma^7 is real, omega^2 = -1, Tr omega = 0; omega anticommutes with every gamma^A and commutes with all 28 generators Sigma^{AB}, so that the eigenspaces of chi are invariant under Spin_0
[PASS] 71-R2  7+1: the linear system X M^{AB} = M^{AB} X (all 28 generators; 7168 integer equations for the 256 entries of X) has rank 254 modulo a prime, so its solutions form a space of dimension at most 2; 1 and omega are two independent solutions; hence the commutant is exactly the span of 1 and omega, of dimension 2, over the real numbers and (the equations being real) also over the complex numbers
[PASS] 71-R3  7+1: the real commutant is spanned by 1 and omega with omega^2 = -1, an algebra isomorphic to the complex numbers; (a + b omega)^2 = a + b omega has only the solutions 0 and 1, so the real commutant contains no projector other than 0 and 1, and the real 16-dimensional representation is IRREDUCIBLE over the real numbers. Over the complex numbers P^L = (1 - i omega)/2 and P^R = (1 + i omega)/2 split it into two 8-dimensional parts, which by check R2 are irreducible and not equivalent
[PASS] 71-R4  7+1: in the basis of the matrix files the generators are not block-diagonal: no choice of real components separates the two chiral parts; the left part consists of complex combinations of the real components
[PASS] 71-R5  7+1: because omega is real, the complex conjugate of 2P^L = 1 - i omega is 2P^R = 1 + i omega, and the generators are real: complex conjugation maps the left part onto the right part and commutes with the action of Spin_0. The right part is the complex conjugate of the left part
[PASS] 71-R6  7+1: the 28 generators remain linearly independent on each part (rank 28 on the left part, 28 on the right part): the whole Lie algebra so(7,1) acts faithfully on each 8-dimensional part; each part is a representation of the one group Spin_0(7,1), not of a subgroup
[PASS] 71-R7  7+1: C^- M^{AB} + (M^{AB})^T C^- = 0 for all generators (C^- is symmetric), and C^- anticommutes with omega: the invariant form psi^T C^- psi' pairs the left part with the right part. The right part is therefore also the dual of the left part
[PASS] 71-R8  7+1: omega = (gamma^0 gamma^1)(gamma^2 gamma^3)(gamma^4 gamma^5)(gamma^6 gamma^7), where the factor with the time axis x4 squares to +1 (a boost) and the other three to -1. omega anticommutes with every gamma^A, so that it maps every vector to its negative; that reverses the direction of time, and therefore omega is not in Spin_0(7,1). The central elements of Spin_0(7,1) are +1 and -1, and -1 acts as -1 on each part: each 8-dimensional part is a faithful representation of Spin_0(7,1)
17 of 17 checks passed
\end{verbatim}

\hypertarget{notation}{%
\section{8. Notation}\label{notation}}

\begin{longtable}[]{@{}
  >{\raggedright\arraybackslash}p{(\columnwidth - 2\tabcolsep) * \real{0.3649}}
  >{\raggedright\arraybackslash}p{(\columnwidth - 2\tabcolsep) * \real{0.6351}}@{}}
\toprule\noalign{}
\begin{minipage}[b]{\linewidth}\raggedright
symbol
\end{minipage} & \begin{minipage}[b]{\linewidth}\raggedright
meaning
\end{minipage} \\
\midrule\noalign{}
\endhead
\bottomrule\noalign{}
\endlastfoot
\(x0,\dots,x7\) & the coordinates: \(x1,x2,x3\) ordinary space, \(x4\)
time, \(x0\) hidden, \(x5,x6,x7\) extra \\
\(\gamma^A\), \(\eta^{AB}\), \(A,B,C\) & the real \(16\times16\) Dirac
matrices; the flat metric; indices \(0,\dots,7\) \\
\(\Sigma^{AB}\), \(M^{AB}\) & \(\tfrac14[\gamma^A,\gamma^B]\), the
generators of the spin group; \(4\Sigma^{AB}\) \\
\(\mathrm{Spin}_0(4,4)\), \(\mathrm{Spin}_0(7,1)\) & the identity
components of the spin groups \\
\(\mathfrak{so}(4,4)\), \(\mathfrak{so}(7,1)\) & their Lie algebras \\
\(SO(4,4)\), \(SL(2,C)\) & the group of real \(8\times8\) matrices
preserving a symmetric form of signature (4,4); the group of complex
\(2\times2\) matrices of determinant 1 \\
\(w_{AB}\) & real parameters of an element
\(\exp(\sum_{A<B}w_{AB}\Sigma^{AB})\) \\
\(\omega\), \(\chi\) & \(\gamma^{0}\cdots\gamma^{7}\); the chirality
matrix, \(\omega\) (4+4) or \(i\omega\) (7+1) \\
\(P^L\), \(P^R\) & \(\tfrac12(1\mp\chi)\), the chiral projectors \\
\(C^-\) & \(\gamma^{0}\gamma^{1}\gamma^{2}\gamma^{3}\) (4+4),
\(\gamma^{4}\omega\) (7+1) \\
\(\tau^A\), \(\bar\tau^A\) & the real \(8\times8\) blocks of the 4+4
Dirac matrices in the files \\
\(\psi\), \(\psi'\) & spinors with sixteen components \\
\(X\), \(E\), \(T\) & a matrix of the commutant; a projector; an
intertwining matrix \\
\(\rho_1\), \(\rho_2\), \(g\) & two representations; a group element (in
Sec. 2 only) \\
\(1_8\) & the \(8\times8\) unit matrix \\
\end{longtable}

\end{document}
